% !TEX root = ../main.tex
\chapter{The Central Bank as a Clearinghouse}\label{ch:clearinghouse}
\begin{paracol}{2}

\begin{sources}
\begin{tabularx}{\linewidth}{@{}l l L@{}}
\toprule
\textbf{Segment} & \textbf{Length} & \textbf{Content}\\
\midrule
\seg{5}{1}{L5.1 Martin Wolf on QE3} & 3:16 & Wolf: QE3 more likely to do too little than too much; the nominal GDP gap.\\
\seg{5}{2}{L5.2 One Big Bank} & 8:20 & A payment system as one big bank: money version and credit version; the payment matrix.\\
\seg{5}{3}{L5.3 Multiple Banks, a challenge} & 3:56 & Many banks: paying in reserves is more constraining.\\
\seg{5}{4}{L5.4 Reading: Charles F.~Dunbar} & 2:05 & Dunbar's \emph{Chapters on the Theory and History of Banking} (1891).\\
\seg{5}{5}{L5.5 Correspondent banking} & 10:24 & Netting, correspondent balances as a swap of IOUs; who expands.\\
\seg{5}{6}{L5.6 Correspondent banking, system network} & 3:55 & The hierarchical network of correspondents up to New York.\\
\seg{5}{7}{L5.7 Clearinghouse, normal operations} & 8:28 & Clearinghouse certificates, multilateral netting, settlement.\\
\seg{5}{8}{L5.8 Clearinghouse, private lender of last resort} & 10:44 & Clearinghouse loan certificates in crises; why the Fed replaced them.\\
\seg{5}{9}{L5.9 Central Bank Clearing} & 4:54 & The Fed as a public clearinghouse: Fedwire, Fed funds, discount loans.\\
\seg{5}{10}{L5.10 Central Bank Cooperation} & 5:41 & Internal and external drain; the major central banks acting as a clearinghouse.\\
\bottomrule
\end{tabularx}

\smallskip
\textbf{Files.} Transcripts \file{materials/transcripts/L05-P*.txt}; Mehrling's notes \mn{5}{1}.
\textbf{Reading.} Charles F.~Dunbar, \emph{Chapters on the Theory and History of Banking} (1891; 2nd ed.\ 1901), the chapter on the
checking system (not among the course files).
\end{sources}

\section*{Motivation}
This lecture opens the first block of the course, \emph{banking as a system of payments}, and \emph{central banking as a clearinghouse}
\ts{5}{2}{0:08}. Starting from payments is unusual, but it builds the intuition for what follows: why the Fed funds market is so central,
what goes on in repo and in the Eurodollar market \ts{5}{2}{0:28}. The guiding idea \ts{5}{2}{0:48}:
\begin{center}
\emph{all the institutions and practices of the monetary system are attempts, worked out over a long time, to make a system of many banks
work as if it were one big bank.}
\end{center}
The lecture goes up the hierarchy: one big bank (\cref{sec:ch-onebank}); many banks (\cref{sec:ch-many}) linked by correspondent balances
(\cref{sec:ch-corr}); the clearinghouse at the top, in normal times and in crises (\cref{sec:ch-chouse,sec:ch-lolr}); the Fed as a public
clearinghouse (\cref{sec:ch-fed}); and the central banks of the world as a clearinghouse one level higher (\cref{sec:ch-coop}).

% ------------------------------------------------------------------
\section{FT: Martin Wolf on QE3}\label{sec:ch-ft}

\begin{casestudy}[FT, September 2012: ``Bernanke makes a historic choice'']
Martin Wolf, the FT's chief economics commentator, writes every Wednesday \ts{5}{1}{0:09}. He argues that critics of the new Fed policy
are right that it will probably fail to work as hoped, but wrong that it will do vast damage: the fear of hyperinflation is misguided, and
the costs of deficient demand are far greater; the Fed ``is in truth more likely to achieve too little of what it seeks than too much''
\ts{5}{1}{0:50}. Mehrling agrees: it is not clear how QE translates into aggregate demand \ts{5}{1}{1:10}. Wolf's chart of nominal GDP
against its pre-crisis trend, the goal that Bernanke and Michael Woodford discussed at Jackson Hole, shows a gap that is widening, not
closing; perhaps because fiscal policy is stuck, too much is being asked of the central bank, in the US and in Europe
\ts{5}{1}{1:50}.
\end{casestudy}\begin{aside}[The Money View blog]
Mehrling and his co-blogger Daniel Neilson posted on QE3 on their blog, together with a short video on Draghi and the Soros plan
(\cref{sec:nh-ft}) \ts{5}{1}{1:30}.
\end{aside}

% ------------------------------------------------------------------
\section{One big bank}\label{sec:ch-onebank}

``One big bank'' can mean two things, following the money/credit distinction of lecture~4 \ts{5}{2}{1:29}.

\paragraph{The money version.} The bank holds reserves (say gold) and owes deposits to people, $\alpha$ and $\beta$. A payment from
$\alpha$ to $\beta$ is a book entry from one account to another; no reserves move at all \ts{5}{2}{2:11}. The limit: when a
deposit is run down to zero, no more purchases are possible, unless one can borrow, bilaterally, from someone who has a deposit
\ts{5}{2}{2:54}. The IMF, with its Special Drawing Rights, is built on this discipline model: a fixed balance sheet, reallocated,
expanded only now and then ``in one big tranche'' \ts{5}{2}{3:14}.

\paragraph{The credit version.} The bank has deposits and also \emph{overdrafts} of people $\gamma$ and $\lambda$, on the asset side: it
uses the fact that no reserves are really needed \ts{5}{2}{3:35}. Now the balance sheet may expand or contract, depending on the
pattern of payments \ts{5}{2}{4:16}.
\begin{taccts}
\tacct[2.4cm]{One big bank (money)}{Reserves}{Deposit $\alpha$\\Deposit $\beta$}\qquad
\tacct[2.4cm]{One big bank (credit)}{Overdraft $\gamma$\\Overdraft $\lambda$}{Deposit $\alpha$\\Deposit $\beta$}
\end{taccts}

\begin{nfigure}
\renewcommand{\arraystretch}{1.35}
\begin{tabular}{c|cc|cc}
from $\backslash$ to & $\alpha$ & $\beta$ & $\gamma$ & $\lambda$\\\hline
$\alpha$ & 0 & 0 & $-$ & $-$\\
$\beta$ & 0 & 0 & $-$ & $-$\\\hline
$\gamma$ & $+$ & $+$ & 0 & 0\\
$\lambda$ & $+$ & $+$ & 0 & 0\\
\end{tabular}
\caption{The payment matrix of the credit one-big-bank \ts{5}{2}{4:37}: change in the size of the balance sheet when the row agent pays the
column agent. $\alpha,\beta$ have deposits, $\gamma,\lambda$ overdrafts.}\label{fig:ch-matrix}
\end{nfigure}

Reading the matrix \ts{5}{2}{4:58}:
\begin{itemize}
  \item paying oneself, or a depositor paying a depositor: no change, exactly as in the money version;
  \item an overdrawn agent paying another overdrawn agent: one overdraft grows, the other shrinks, total unchanged;
  \item a depositor ($\alpha$) paying an overdrawn agent ($\gamma$): $\alpha$'s deposit falls and $\gamma$'s overdraft is partly repaid:
        the balance sheet \emph{contracts} \ts{5}{2}{6:21};
  \item an overdrawn agent paying a depositor: the overdraft grows and the deposit grows: the balance sheet \emph{expands} \ts{5}{2}{6:41}.
\end{itemize}
The aim of a payment system is that whenever there is a willing buyer and a willing seller, the payment system does not stand in the way.
The credit system can do everything the money system does and more; the actual US domestic payment system has this credit character, and
the Federal Reserve System is an example of it \ts{5}{2}{7:22}.

\begin{secondary}[What is the quantity of money?]
In the money version it is natural to call the deposits the quantity of money. With overdrafts the measure becomes ambiguous; the notes list three
candidates \mn{5}{1}:
\begin{enumerate}
  \item deposits, which counts only the positive balances;
  \item deposits minus overdrafts;
  \item deposits minus credit limits (the unused borrowing power counts as means of payment too).
\end{enumerate}
This ambiguity led Fischer Black to propose not trying to measure money at all, in ``Banking in a World without Money'' (lecture~\ref{ch:threeviews})
\mn{5}{2}.
\end{secondary}

\begin{definition}[New concepts: overdraft]
An \emph{overdraft} is a negative balance on a deposit account, i.e.\ a loan from the bank created automatically when the holder pays more
than he has (one ``draws over'' the balance). It is an asset of the bank.
\end{definition}

% ------------------------------------------------------------------
\section{Many banks}\label{sec:ch-many}

There is no one big bank. In the United States there are thousands of banks, fewer than the six or seven thousand of Young's time, some very
big (Citibank, Bank of America), many small or regional; in some countries, such as the UK or Canada, there are a handful of big banks
\ts{5}{3}{0:07}. The problem is how to knit them into one country-wide system in which anybody can pay anybody
\ts{5}{3}{0:48}.

Take bank A with depositors $\alpha,\beta$ and bank B with depositors $\gamma,\lambda$, all with positive deposits; the credit element now
is \emph{interbank} credit. This is the wholesale money market, behind the scenes, and its elasticity is what provides elasticity in the retail
market \ts{5}{3}{1:09}. One could pay in reserves: when $\alpha$ pays $\gamma$ 10, bank A sends 10 of reserves to
bank B \ts{5}{3}{2:31}:
\begin{taccts}
\tacct[2.2cm]{Bank A}{$-$Reserves 10}{$-$Deposit $\alpha$ 10}\qquad
\tacct[2.2cm]{Bank B}{$+$Reserves 10}{$+$Deposit $\gamma$ 10}
\end{taccts}
This is \emph{more} constraining than one big bank: not only must $\alpha$ have a deposit, bank A must have reserves. So banks do not
generally do this within a country \ts{5}{3}{3:12}.

% ------------------------------------------------------------------
\section{Reading: Charles F.~Dunbar}\label{sec:ch-dunbar}

\begin{reading}[Dunbar, \emph{Chapters on the Theory and History of Banking} (1891)]
This was the textbook used at Harvard College; the course reads its chapter on how the checking system works in the United States
\ts{5}{4}{0:08}. It explains a payment system without a central bank (there was no Fed in 1891), which helps to understand why one was
eventually created; and it is ``crystal clear'': no four-colour diagrams, no workbook, no slides \ts{5}{4}{0:29}. Dunbar was
the first professor of money and banking at Harvard, a former journalist who had written on banking for a Boston paper; the second edition
was enlarged by O.~M.~W.~Sprague, of the Harvard Business School, a good friend of Allyn Young \ts{5}{4}{1:09}.
\end{reading}\begin{history}[Dunbar and Sprague]
Charles Franklin Dunbar (1830--1900) edited the \emph{Boston Daily Advertiser} before becoming, in 1871, the first professor of political economy
at Harvard; he founded the \emph{Quarterly Journal of Economics} in 1886. The second edition of his \emph{Chapters} (1901) was edited by Oliver
M.~W.~Sprague (1873--1953), later the author of \emph{History of Crises under the National Banking System} (1910) for the National Monetary
Commission. In class the paper is called ``the Boston Globe'' \ts{5}{4}{1:09}.\par
\emph{References:} F.~W.~Taussig, ``Charles Franklin Dunbar'', \emph{QJE} 14 (1900); Dunbar, \emph{Chapters on the Theory and History of
Banking}, 2nd ed., Putnam, 1901.
\end{history}

% ------------------------------------------------------------------
\section{Correspondent banking}\label{sec:ch-corr}

\paragraph{Collecting and netting.} During the day bank A receives checks, orders to pay. It does not pay them one by one: it collects them
and sorts them by \emph{bank} (the routing number at the bottom of a check), not by customer; which account at bank B gets credited is bank B's
problem \ts{5}{5}{0:27}. Checks between customers of bank A itself need no reserve flows; for each other bank it
sums what it owes (``due to B'') and what it is owed (``due from B''), and at the end of the day computes the net, which may be positive or
negative \ts{5}{5}{1:29}. Netting already saves moving reserves back and forth across the street every
day \ts{5}{5}{3:43}.

\paragraph{Correspondent balances: a swap of IOUs.} Banks that do business regularly can do better: bank A holds a deposit at B in its own
name (a \emph{banker's balance}), and B holds one at A. Instead of paying for these deposits with reserves, they swap them: each says to the
other ``I owe you \$1000''. Both balance sheets expand, deposits ``created from thin air'' \ts{5}{5}{4:04}.
\begin{taccts}
\tacct[2.3cm]{Bank A}{$+$Deposit at B}{$+$Deposit of B}\qquad
\tacct[2.3cm]{Bank B}{$+$Deposit at A}{$+$Deposit of A}
\end{taccts}
It looks like a free lunch that cannot do anything, but it does: now net payments can be made without moving reserves
\ts{5}{5}{5:30}. For $\alpha$ (at A) to pay $\lambda$ (at B) 10, there are two ways \ts{5}{5}{5:51}--\ts{5}{5}{8:16}:
\begin{enumerate}
  \item A draws down its deposit at B: A's balance sheet \emph{shrinks}, B's only changes owner (B swaps an IOU to A for an IOU to $\lambda$);
  \item B increases its deposit at A (A's liabilities to B grow): A's balance sheet does not change in size, B's \emph{expands}.
\end{enumerate}
\begin{taccts}
\textit{(1)}\quad\tacct[2.2cm]{Bank A}{$-$Deposit at B 10}{$-$Deposit $\alpha$ 10}\quad
\tacct[2.2cm]{Bank B}{}{$-$Deposit of A 10\\$+$Deposit $\lambda$ 10}

\medskip
\textit{(2)}\quad\tacct[2.2cm]{Bank A}{}{$-$Deposit $\alpha$ 10\\$+$Deposit of B 10}\quad
\tacct[2.2cm]{Bank B}{$+$Deposit at A 10}{$+$Deposit $\lambda$ 10}
\end{taccts}
Either way it is credit: either the debtor bank's balance sheet shrinks or the creditor bank's expands \ts{5}{5}{8:16}. In the first case total
deposits, counting interbank balances, fall; in the second they rise \mn{5}{3}. The invention of correspondent banking moves payments \emph{between
banks} from a money system to a credit system: they are no longer constrained by the quantity of gold, only by bilateral credit limits \mn{5}{3}.

\begin{definition}[New concepts: correspondent balances]
A \emph{correspondent balance} (banker's balance) is a deposit that one bank holds at another to settle payments between them. The bank that
holds the account is the \emph{respondent}, the bank that keeps it the \emph{correspondent} (the two ``correspond'', i.e.\ exchange letters and
orders, in the old sense of the word).
\end{definition}

\paragraph{Which one gets picked? Hierarchy.} In practice the big city bank says to the small country bank: I have no reason to hold deposits
with you; you hold a deposit with me, paid in reserves up front, or I lend it to you for a fee, and my deposits are the clearing balances
\ts{5}{5}{8:37}. This is discipline for the debtor bank: to join the payment system a country bank must keep balances at an urban
bank; if it runs out there is no guarantee the city bank will expand to help it, and it may charge extra. ``This is the hierarchy of money
showing up in the payment system'' \ts{5}{5}{9:17}. Correspondent balances make the system work much more like the ideal one big
bank on a credit basis; the retail customer sees none of this \ts{5}{5}{9:58}. The notes: the more central bank always takes deposits from the
less central; in the US country banks held deposits at city banks, and those reserve cities survive as the cities of the twelve Federal Reserve Banks
\mn{5}{3}.

\paragraph{The network.} Before the central bank, money-centre banks (Cleveland, Chicago, \dots) had correspondent relations with banks in their
hinterland; payments within a region just moved deposits at the centre bank \ts{5}{6}{0:07}. Payments between regions need a payment
one layer up, between centre banks; in practice this was not done bilaterally but hierarchically, through a few New York banks at which all the
money centres held accounts, with further regional layers \ts{5}{6}{1:10}. The whole structure nets as
much as it can, so that reserves do not have to flow, and builds credit relations locally and higher up, balancing elasticity and discipline from
the smallest bank ``in the middle of nowhere'' up to J.~P.~Morgan in New York \ts{5}{6}{2:53}. At the top, the biggest New York banks
must clear with each other: that is where the clearinghouse comes in \ts{5}{6}{3:34}.

\begin{nfigure}
\begin{tikzpicture}[font=\scriptsize,
  ny/.style={draw,thick,fill=keycol!20,rounded corners=2pt,minimum width=1.6cm,minimum height=0.6cm},
  mc/.style={draw,fill=defcol!15,rounded corners=2pt,minimum width=1.3cm,minimum height=0.5cm},
  cb/.style={draw,fill=fincol!15,circle,inner sep=1.6pt}]
  \node[ny] (ny) at (0,3) {New York bank};
  \foreach \x/\n [count=\i] in {-4.2/Cleveland,0/Chicago,4.2/St.\ Louis} {
    \node[mc] (m\i) at (\x,1.5) {\n};
    \draw[semithick] (m\i) -- (ny);
    \foreach \d in {-1.2,-0.6,0,0.6,1.2} {
      \draw[thin] (\x+\d,0) -- (m\i);
      \node[cb] at (\x+\d,0) {};}}
  \node[right,font=\scriptsize] at (5.6,0) {country banks};
  \node[right,font=\scriptsize] at (5.6,1.5) {money-centre banks};
  \node[right,font=\scriptsize] at (5.6,3) {New York};
\end{tikzpicture}
\caption{The correspondent network before the Fed \ts{5}{6}{0:07}: country banks hold balances at their money-centre bank, money-centre
banks at a New York bank; a payment between two country banks in different regions moves balances at every level up to the common correspondent.}
\label{fig:ch-network}
\end{nfigure}

% ------------------------------------------------------------------
\section{The clearinghouse in normal times}\label{sec:ch-chouse}

The New York Clearing House Association (NYCHA) is a mutual organisation of the big New York banks, ``first among equals'' J.~P.~Morgan
\ts{5}{7}{0:08}. It was set up on the discipline model: every member deposits gold or currency, and the clearinghouse issues
\emph{clearinghouse certificates} to its credit, physical notes standing for the reserves \ts{5}{7}{0:55}.
\begin{taccts}
\tacct[2.6cm]{Clearinghouse}{Gold (reserves)}{Clearinghouse certificates (to A, B, \dots)}
\end{taccts}

\paragraph{During the day: credit.} The members agree to treat every amount due to or from another member as due to or from the clearinghouse
itself \ts{5}{7}{1:58}. Each bank then has a single net position against the clearinghouse, smaller than the sum of its bilateral nets,
because positives and negatives offset: \emph{multilateral netting} \ts{5}{7}{2:40}. In Dunbar's description these are
actual checks: at the end of the day the banks add them up and meet at the clearinghouse, and those in net deficit hand clearinghouse certificates to
those in net surplus; then they go home \ts{5}{7}{3:42}.

\begin{intuition}[Bilateral versus multilateral netting]
With $n$ banks and gross obligations $x_{ij}$ (from $i$ to $j$), bilateral netting settles $|x_{ij}-x_{ji}|$ for each pair, while multilateral
netting settles only each bank's net position $\sum_j(x_{ij}-x_{ji})$. The sum of the multilateral nets is never larger than the sum of the
bilateral ones (triangle inequality), and is usually much smaller: a ring of payments $A\to B\to C\to A$ of equal size nets to zero.
\end{intuition}

``There is a credit expansion during the day that creates elasticity, and then there's a settlement at the end of the day'': elasticity and
discipline at the very highest level of the system \ts{5}{7}{4:23}. The members know they have this elasticity and can sell it to the banks
below them, and so on down: the whole system operates as a credit system because the clearinghouse does, during the day. Sometimes settlement is
every three days, which is why the length of the \emph{float} matters \ts{5}{7}{5:05}.

\paragraph{Failing to settle.} A bank that has run down its certificates over several days of deficits has three options
\ts{5}{7}{5:48}:
\begin{enumerate}
  \item pay in clearinghouse certificates, if it has them;
  \item borrow certificates bilaterally from another member, e.g.\ from J.~P.~Morgan, or from a member in net surplus;
  \item if nobody will lend, because it is thought insolvent, default: the clearinghouse's gold goes to pay the creditors pro rata, the losses are
        shared among all the members, and the defaulter is expelled, ``banished to the outer ring of hell''.
\end{enumerate}

\begin{proposition}[Joint credit]
The credit of a group of banks that jointly and severally guarantee an obligation is better than the credit of any individual bank: clearinghouse
credit is better than the bilateral credit of its members \ts{5}{7}{7:32}.
\end{proposition}

\begin{definition}[New concepts: clearing, settlement, netting, float]
\emph{Clearing} is the process of collecting and offsetting payment orders between banks; \emph{settlement} is the final transfer of the means of
payment (reserves) that discharges the net obligations. \emph{Netting} replaces gross obligations by net ones. The \emph{float} is the credit that
exists between the moment a payment is credited to the payee and the moment it is settled.
\end{definition}

% ------------------------------------------------------------------
\section{The clearinghouse as private lender of last resort}\label{sec:ch-lolr}

In times of stress all the members need to make payments, to their correspondents or abroad, in gold; they all want to withdraw their accounts at
the clearinghouse, and even the biggest banks can run out of certificates: the country banks hold their reserves in New York, and the New York banks
do not have enough gold to cover them \ts{5}{8}{0:08}. Then the clearinghouse creates money, not by bilateral borrowing but
by joint borrowing and lending \ts{5}{8}{0:50}: it lends to bank A (at 6\% in the example) by printing a \emph{clearinghouse loan
certificate}, a liability of the clearinghouse, guaranteed by all its members, against whatever collateral A can pledge, for example government
securities \ts{5}{8}{1:32}.
\begin{taccts}
\tacct[2.6cm]{Clearinghouse}{$+$Loan to A (6\%)}{$+$Loan certificates}\qquad
\tacct[2.6cm]{Bank A}{$+$Loan certificates}{$+$Loan from clearinghouse}
\end{taccts}
Money has been created at the apex of the system: more of the best money, which circulates among banks as if it were gold. The certificates say
``loan certificate'', but they are treated exactly like the ordinary certificates, since both are liabilities of the clearinghouse
\ts{5}{8}{2:38}. This is a lender of last resort operation, managed by the clearinghouse; the New York
clearinghouse still exists, still clears for its members, and clears much of the Eurodollar market \ts{5}{8}{3:40}.

\paragraph{Bagehot's rule in private hands.} A committee decides \ts{5}{8}{4:21}. Six percent in 1907 was a high rate, ``lend freely at a high rate
against good security'': high, so that the borrower wants to repay. So high, in fact, that banks liked to keep these certificates, the safest asset in
the system at 6\%, and after the crisis the clearinghouse had trouble getting them back \ts{5}{8}{4:42}. Good security: the
bankers know each other and would not take ``garbage''; a bank without acceptable collateral had to borrow some from another bank first, unlike the Fed,
which can decide to let it pass \ts{5}{8}{5:43}.

\paragraph{Playing favourites.} Young describes how, in crises, New York banks refused to pay the deposits of country banks in reserves, offering loan
certificates, take it or leave it, and kept their gold to pay abroad, where the Bank of England would accept nothing but proper international money
\ts{5}{8}{6:46}. It was never clear that the certificates were legal, since they were used as money (unlike the bank notes
backed by government bonds of \cref{sec:st-nbsorigin}) \ts{5}{8}{7:47}; as long as they stayed among bankers nobody noticed, but sometimes
they circulated. After 1907 a law made the practice legal, with rules on acceptable collateral \ts{5}{8}{8:29}. Some argue that this is all
that should have been done, and there is some truth in it: it was a decentralised lender of last resort, discipline in good times and elasticity in
crises \ts{5}{8}{9:11}. But the New York bankers were playing favourites; the banks of the hinterland, refused payment when
everybody needed it, sometimes went bankrupt to save the New York banks, and pressed for democratic control. That is how the Fed came about: ``one way to
understand the Fed is that it is basically a clearinghouse'' \ts{5}{8}{10:13}.

\begin{coursenote}[History of the NYCHA]
The New York Clearing House Association was founded in 1853, long before Morgan's prominence; Morgan's role as ``first among equals'' refers to the
crises of the 1890s and 1907. Clearinghouse loan certificates were first issued in New York in 1857 and 1860 and in most crises thereafter (1873, 1884,
1890, 1893, 1907). The law Mehrling refers to is presumably the Aldrich--Vreeland Act (1908), which authorised groups of national banks
(``national currency associations'') to issue emergency currency against commercial paper and other securities; it was used once, in 1914. Today the
association's successor, The Clearing House, operates CHIPS.
\end{coursenote}\begin{history}[The panic of 1907]
In October 1907 the run on the Knickerbocker Trust Company spread to the trust companies, which were not members of the New York Clearing House. J.~P.~Morgan
assembled bankers in his library to organise loans to the weaker institutions, and the clearinghouses of New York and of many other cities issued loan
certificates; currency went to a premium over certified checks, i.e.\ par between currency and deposits was broken.\par
\emph{References:} Bruner and Carr, \emph{The Panic of 1907}, 2007; Sprague, \emph{History of Crises under the National Banking System}, 1910.
\end{history}

% ------------------------------------------------------------------
\section{The Fed as a public clearinghouse}\label{sec:ch-fed}

The Federal Reserve, when on the gold standard, held gold and issued Federal Reserve notes and bank reserves (deposits of banks at the Fed)
\ts{5}{9}{0:07}. It can run intraday credit on its own books: \emph{Fedwire}, now electronic \ts{5}{9}{0:49}. The private clearing system
of the New York banks is today called \emph{CHIPS}. CHIPS clears first, then Fedwire, and Fedwire is ``real money'', deposits at the Fed, the ultimate
settlement asset \ts{5}{9}{1:12}.
\begin{taccts}
\tacct[2.6cm]{Federal Reserve}{Gold\\Discount loans}{Federal Reserve notes\\Bank reserves}
\end{taccts}
Overdrafts and balances build up during the day. At the end of the day a bank without a positive balance at the Fed must \ts{5}{9}{1:33}:
\begin{enumerate}
  \item borrow from another bank that has a surplus: the \emph{Fed funds market} (next lecture);
  \item or borrow from the Fed, which makes \emph{discount loans} by expanding its own balance sheet, if the banks cannot find each other or B is
        unwilling: nobody forces them \ts{5}{9}{2:16};
  \item or default.
\end{enumerate}
Discount loans to troubled banks appear on the Fed's balance sheet and are small: with thousands of banks, one fails every week or so, and whole
divisions deal with them; they matter for the aggregate only in system-wide stress \ts{5}{9}{2:58}. Then the Fed does what the
clearinghouse did, but legally, as monetary policy: it expands both sides of its balance sheet to create more of the best money. Bernanke goes further than
any clearinghouse would: he does not only lend to banks, he buys mortgage-backed securities from whoever holds them, and lends to broker-dealers; the
clearinghouse was a club, the Fed has more members and in a crisis deals with the whole economy \ts{5}{9}{3:18}.
Private and public clearinghouse differ in membership, rules and priorities, but the mechanics are very much the same \ts{5}{9}{4:41}.

\begin{intuition}[One step beyond the clearinghouse]
In the notes, central banking is ``nothing more than one step beyond the clearinghouse'': it regularizes and strengthens the clearinghouse and
erases the difference between clearinghouse certificates (backed by gold) and clearinghouse \emph{loan} certificates (backed by members' loans), both
becoming deposits or currency of the central bank \mn{5}{5}. A clearinghouse can be broken by any single member who insists on redeeming its
certificates in gold; a central bank cannot be broken by internal forces, because its liabilities \emph{are} the member banks' reserves. All payments
inside the system net out on its books: it makes the system of banks work as if it were a single bank. What remains is managing relations with the
rest of the world, internal and external drains, and that is the origin of the ``art of central banking'' \mn{5}{5}.
\end{intuition}

\begin{definition}[New concepts: Fedwire, CHIPS]
\emph{Fedwire} is the Federal Reserve's real-time gross settlement system, which transfers reserve balances between banks' accounts at the Fed.
\emph{CHIPS} (Clearing House Interbank Payments System), operated by The Clearing House since 1970, is the private system on which large-value, mostly
international dollar payments are netted among member banks and then settled through their accounts at the Fed.
\end{definition}

% ------------------------------------------------------------------
\section{Central bank cooperation}\label{sec:ch-coop}

What can kill a clearinghouse is members withdrawing their gold because they need it \ts{5}{10}{0:08}. Why would they? Distinguish
\ts{5}{10}{0:28}:

\begin{definition}[New concepts: internal and external drain]
An \emph{internal drain} is a loss of reserves to the domestic public, which wants cash; an \emph{external drain} is a loss of reserves to foreigners,
who want the international money (gold). An internal drain can be met by expanding the balance sheet at the top of the national hierarchy; an
external drain comes from a level higher than the national system, and can kill a clearinghouse or a central bank.
\end{definition}

The same holds for central banks: if foreign central banks bring Federal Reserve notes and ask for gold, the Fed can run out, because its notes are
not fully covered \ts{5}{10}{1:09}. Then it \emph{suspends} payments: ``I know these were promises to pay gold, I changed my mind; I will pay
some time later'': it abandons the fixed mint parity, and all central banks have done it, the Bank of England periodically
\ts{5}{10}{1:51}. An internal crisis can be handled by expanding the balance sheet; when the whole country loses reserves, the only way to
get elasticity is to eliminate the discipline by going off gold \ts{5}{10}{2:32}, unless all the central banks agree to act as a clearinghouse
and expand together, one layer higher \ts{5}{10}{2:53}.

\paragraph{The world clearinghouse.} That is what is happening now: Draghi and Bernanke announce QE; the Bank of England and the Swiss National Bank
already do it; the Bank of Japan was in the FT the day before. ``There are only five central banks that matter'', and whether or not they met in a room
like the New York bankers, they act in concert: all expanding at the same time, so there is no external drain, ``the Martians are not yet doing business
with us'' \ts{5}{10}{3:15}. They do what the IMF cannot, since the IMF needs all its members to agree to expand
SDRs; the ``C5'' can do it privately, leaving the rest of the world out, like the New York bankers. Expect complaints from the rest of the world that the
big central banks take care of each other and not of them \ts{5}{10}{4:37}.

\begin{nfigure}
\begin{tikzpicture}[font=\scriptsize,
  l/.style={draw,rounded corners=2pt,align=center,minimum width=3.3cm,minimum height=0.8cm}]
  \node[l,fill=keycol!15] (w) at (0,3.6) {central banks acting together\\(external drain disappears)};
  \node[l,fill=defcol!15] (c) at (0,2.4) {central bank / clearinghouse\\(internal drain, lender of last resort)};
  \node[l,fill=intcol!15] (m) at (0,1.2) {money-centre banks\\(correspondent balances)};
  \node[l,fill=fincol!15] (b) at (0,0) {banks and their customers\\(deposits, overdrafts)};
  \foreach \u/\d in {w/c,c/m,m/b} \draw[-{Stealth[length=4pt]}] (\u) -- (\d);
  \node[right,align=left] at (2.4,1.8) {each level: credit during the day (elasticity),\\settlement in the asset of the level above (discipline)};
\end{tikzpicture}
\caption{Clearing all the way up: the same mechanism of intraday credit and final settlement repeats at each level of the hierarchy, from bank customers to
the central banks of the world.}\label{fig:ch-levels}
\end{nfigure}

% ------------------------------------------------------------------
\flushside
\section*{Key formulas and ideas}
\begin{keyformulas}
\begin{tabularx}{\linewidth}{@{}l L@{}}
One big bank & money version: payments are book entries, balance sheet fixed; credit version: depositor$\to$overdrawn contracts, overdrawn$\to$depositor expands\\
Many banks & paying in reserves is more constraining than one big bank\\
Correspondent balances & a swap of IOUs; a payment either shrinks the debtor bank or expands the creditor bank; hierarchy decides which\\
Clearinghouse & certificates against gold; intraday multilateral netting (elasticity); end-of-day settlement in certificates (discipline)\\
Failing to settle & pay, borrow from a member, or default (losses shared pro rata)\\
Loan certificates & clearinghouse lends its own liabilities against collateral at a high rate: private lender of last resort\\
Fed & public clearinghouse: Fedwire, Fed funds, discount loans\\
Drains & internal: met by expansion; external: suspension, unless central banks expand together\\
\end{tabularx}
\end{keyformulas}

\section*{Exercises}

\begin{exercise}[The payment matrix]
In the credit one-big-bank, $\alpha$ and $\beta$ have deposits of 50 each, $\gamma$ and $\lambda$ overdrafts of 30 each. Compute the size of the balance sheet
after each of the following payments, applied in sequence: $\gamma$ pays $\alpha$ 20; $\alpha$ pays $\lambda$ 40; $\lambda$ pays $\gamma$ 10.
\end{exercise}

\begin{exercise}[Netting]
Four banks have the following gross obligations at the end of the day (row pays column): A$\to$B 30, B$\to$C 40, C$\to$A 35, A$\to$C 10, D$\to$A 5, B$\to$D 15.
\begin{enumerate}[(a)]
  \item Compute the total gross payments, the total of bilateral nets, and the total of multilateral nets.
  \item Which banks must deliver certificates to the clearinghouse, and how many?
\end{enumerate}
\end{exercise}

\begin{exercise}[Correspondent balances]
Country bank C holds a balance of 100 at city bank N. Customers of C pay customers of N 30 more than they receive.
\begin{enumerate}[(a)]
  \item Write the T-accounts when the net is settled by drawing down C's balance.
  \item The balance is now too low; N offers to lend C 50 by crediting its account. Write the T-accounts and explain why this is ``discipline for the debtor
        bank''.
\end{enumerate}
\end{exercise}

\begin{exercise}[Loan certificates]
Bank A owes the clearinghouse 20 at settlement and has no certificates. The clearinghouse issues 20 of loan certificates at 6\% against A's government bonds.
Write the T-accounts of the clearinghouse, A, and the creditor bank B. What happens a week later when A repays? Why might B prefer not to be repaid?
\end{exercise}
